\documentclass{article}

\usepackage{booktabs}
\usepackage{tabularx}
\usepackage{float}

\input{../Comments}
%% Common Parts

\newcommand{\progname}{Project Proxi} % PUT YOUR PROGRAM NAME HERE
\newcommand{\authname}{Team \#23, Project Proxi
\\ Savinay Chhabra
\\ Amanbeer Singh Minhas
\\ Gourob Podder
\\ Ajay Singh Grewal} % AUTHOR NAMES                  

\usepackage{hyperref}
    \hypersetup{colorlinks=true, linkcolor=blue, citecolor=blue, filecolor=blue,
                urlcolor=blue, unicode=false}
    \urlstyle{same}
                                


\title{User Guide\\\progname}

\author{\authname}

\date{}

\begin{document}

\setlength{\parindent}{0pt}

\begin{table}[hp]
\caption{Revision History} \label{TblRevisionHistory}
\begin{tabularx}{\textwidth}{llX}
\toprule
\textbf{Date} & \textbf{Developer(s)} & \textbf{Change}\\
\midrule
April 5, 2026 & Savinay & Initial User Guide release with interface startup, TUI/GUI settings, MCP setup, and Discord integration guidance\\
April 6, 2026 & Savinay & Added \texttt{uv run proxi setup} installation flow and documented the GUI voice controls\\
\bottomrule
\end{tabularx}
\end{table}

\newpage

\maketitle

\newpage

\section{Overview}

This guide describes the recommended workflow for Proxi:
\begin{enumerate}
	\item Start exactly one gateway process
	\item Choose either the TUI or the GUI as your interface
	\item If you use the GUI, configure settings (API keys, Integrations, and cron jobs)
\end{enumerate}

\section{Installation}

From the repository root, install dependencies and initialize the local Proxi setup with:

\begin{verbatim}
uv sync
uv run proxi setup
\end{verbatim}

This one-time setup flow:

\begin{itemize}
	\item Installs Python dependencies with \texttt{uv}
	\item Installs the Node dependencies needed by the TUI, GUI, and Discord relay
	\item Initializes or verifies the local API key database in \texttt{config/api\_keys.db}
\end{itemize}

If you ever need to repair a local environment, rerun \texttt{uv run proxi setup}. It is safe to run more than once.

\section{Starting the Interface}

This section describes the required startup order and your three interface options.

\subsection{Step 1 (Always): Start the Gateway}
Run:

\begin{verbatim}
uv run proxi gateway start
\end{verbatim}

\subsection{Step 2: Choose an Interface}

\subsubsection{Option A: TUI}
\begin{quote}
Run:

\begin{verbatim}
uv run proxi
\end{verbatim}
\end{quote}

\subsubsection{Option B: GUI (React Frontend)}
\begin{quote}
Run:

\begin{verbatim}
uv run proxi frontend
\end{verbatim}
\end{quote}

\subsubsection{Option C: Discord}
\begin{quote}
Run:

\begin{verbatim}
uv run proxi discord
\end{verbatim}

Then send commands via connected Discord bot.
\end{quote}

\section{TUI Interface and Settings}

This section explains only TUI usage and TUI-specific settings/actions.

\subsection{Launch and Connect}

Start the TUI:
\begin{verbatim}
uv run proxi
\end{verbatim}

If not already running, start gateway first:
\begin{verbatim}
uv run proxi gateway start
\end{verbatim}

\subsection{Opening the Command Palette}

In the TUI input box, type:
\begin{verbatim}
/
\end{verbatim}
bun install
bun run start
\begin{itemize}
	\item Type to filter commands
	\item Use up/down arrows to select
	\item Press Enter to execute
	\item Press Esc to close the palette
\end{itemize}

\subsection{Complete Slash Command Reference}

\begin{table}[H]
\centering
\caption{TUI Slash Commands} \label{TblTUISlashCommands}
\begin{tabularx}{\textwidth}{lX}
\toprule
\textbf{Command} & \textbf{What it does} \\
\midrule
\texttt{/agent} & Switch active agent or create a new one \\
\texttt{/switch-agent} & Alias for \texttt{/agent} \\
\texttt{/delete} & Delete current agent (gateway configuration and agent workspace) \\
\texttt{/mcps} & Open MCP management to enable/disable integrations \\
\texttt{/work-dir} & View or change the working directory used by coding/file tools \\
\texttt{/clear} & Clear chat UI and current session history for a fresh session \\
\texttt{/plan} & Open/view current session plan \\
\texttt{/todos} & Open/view current session todos \\
\texttt{/usage} & Show context-window and turn usage statistics \\
\texttt{/help} & Print command reference in chat \\
\texttt{/exit} & Exit the TUI \\
\bottomrule
\end{tabularx}
\end{table}

\subsection{When to Use Which Command}

\begin{itemize}
	\item Use \texttt{/agent} when switching between workflows (for example, planning vs coding)
	\item Use \texttt{/mcps} after adding keys to turn integrations on/off
	\item Use \texttt{/work-dir} before coding tasks to constrain file/tool scope
	\item Use \texttt{/clear} when you want a clean context for a new task
	\item Use \texttt{/plan} and \texttt{/todos} to inspect execution state
	\item Use \texttt{/usage} to monitor context usage and avoid overflow
\end{itemize}

\subsection{TUI Settings Summary}

In TUI mode, operational settings are primarily controlled by slash commands:
\begin{itemize}
    \item Agent setting: \texttt{/agent}
    \item MCP enable/disable: \texttt{/mcps}
    \item Working directory scope: \texttt{/work-dir}
    \item Session reset: \texttt{/clear}
\end{itemize}

\section{GUI Interface and Settings}

Open \textbf{Settings} from the GUI top bar. The settings modal is organized into these sections:

\begin{enumerate}
	\item General
	\item Agent
	\item Profile
	\item API Keys
	\item Integrations
	\item Cron Jobs
	\item Webhooks
	\item Voice
\end{enumerate}

\subsection{General}

Use this section for quick session and model controls.

\subsubsection{What You Can Configure}
\begin{itemize}
	\item \textbf{Clear History}: clears current session history
	\item \textbf{Reasoning Effort}: set reasoning depth (minimal, low, medium, high)
	\item \textbf{Compact Focus Hint}: optional hint used when compacting context
	\item \textbf{Provider and Model}: select active LLM provider/model for new turns
\end{itemize}

\subsubsection{Suggested Usage}
\begin{itemize}
	\item Use \textbf{Clear History} before a new task domain
	\item Use \textbf{Reasoning Effort} to balance speed vs depth
	\item Use \textbf{Compact} when context gets large
	\item Use \textbf{Switch LLM Provider} after selecting provider/model
\end{itemize}

\subsection{Agent}

Use this section to control which agent receives your prompts.

\subsubsection{Fields and Actions}
\begin{itemize}
	\item \textbf{Current Agent}: read-only indicator of active agent ID
	\item \textbf{Agent Selector}: choose another configured agent
	\item \textbf{Switch Agent}: apply selected agent
\end{itemize}

\subsubsection{Recommended Workflow}
\begin{enumerate}
	\item Select the agent that matches your use case (for example planning vs coding)
	\item Switch agent before starting a new workflow
	\item Verify in chat that responses come from the expected persona/behavior
\end{enumerate}

\subsection{Profile}

Use this section to personalize assistant behavior.

\subsubsection{Important Fields}
\begin{itemize}
	\item Name and Occupation: style/context personalization
	\item Location and Timezone: local-time suggestions and schedule accuracy
	\item Email and Email Signature: outbound drafting quality
	\item Demographics: optional context used by selected prompts
\end{itemize}

\subsubsection{Best Practices}
\begin{itemize}
	\item Keep timezone accurate if you use calendar/cron automations
	\item Keep email signature short and reusable
	\item Update profile when switching environments (school/work/personal)
\end{itemize}

\subsection{API Keys}

Use this section to add integration credentials and model-provider keys.

\subsubsection{MCP Key Mapping Table}
Use the following key names exactly as written (do not include \texttt{=} when entering key names in the GUI).

\begin{table}[H]
\centering
\caption{MCP and Provider Key Mapping} \label{TblMCPKeyMapping}
\begin{tabularx}{\textwidth}{lX}
\toprule
\textbf{MCP / Provider} & \textbf{Key Name(s)} \\
\midrule
LLM Provider (OpenAI) & \texttt{OPENAI\_API\_KEY} \\
Google Calendar & \texttt{GMAIL\_CLIENT\_ID}, \texttt{GMAIL\_CLIENT\_SECRET} \\
Gmail & \texttt{GMAIL\_CLIENT\_ID}, \texttt{GMAIL\_CLIENT\_SECRET} \\
Notion & \texttt{NOTION\_API\_KEY}, \texttt{NOTION\_PARENT\_PAGE\_ID} \\
Spotify & \texttt{SPOTIFY\_CLIENT\_ID}, \texttt{SPOTIFY\_CLIENT\_SECRET} \\
Discord relay/config (if used) & \texttt{DISCORD\_BOT\_TOKEN}, \texttt{PROXI\_DISCORD\_COMMAND\_PREFIX} \\
Obsidian & No API key required (vault/path configuration) \\
Weather & Usually no API key required in this setup \\
\bottomrule
\end{tabularx}
\end{table}

\subsubsection{Notes}
\begin{itemize}
	\item \texttt{OPENAI\_API\_KEY} is the primary LLM key in this setup
	\item Keep secrets private and do not commit real values to version control
	\item After adding keys, refresh and verify they are listed
\end{itemize}

\subsection{Integrations}

Use this section to enable or disable integrations after key setup.

\subsubsection{General Enablement Steps}
\begin{enumerate}
	\item Open \textbf{Settings \textgreater\ Integrations}
	\item Toggle required integrations to enabled
	\item Save/refresh
	\item Run one test prompt per integration
\end{enumerate}

\subsection{Cron Jobs}

Use this section to schedule recurring prompts.

\subsubsection{Core Fields}
\begin{itemize}
	\item Source ID
	\item Target Agent
	\item Target Session (optional)
	\item Priority
	\item Prompt
	\item Schedule
\end{itemize}

\subsubsection{Easy Schedule (Recommended)}
If you do not want to write cron patterns manually:
\begin{itemize}
	\item Set \textbf{Every} (number)
	\item Set \textbf{Unit} (Minute, Day, Week, Monthly)
	\item Let the UI generate cron expression automatically
\end{itemize}

Examples:
\begin{itemize}
	\item Every 30 Minute
	\item Every 1 Day
	\item Every 1 Week
\end{itemize}

\subsubsection{Manual Cron (Advanced)}
Example expression:
\begin{verbatim}
*/30 * * * *
\end{verbatim}

Syntax meaning (5-field cron format):
\begin{itemize}
	\item Field 1 = minute
	\item Field 2 = hour
	\item Field 3 = day of month
	\item Field 4 = month
	\item Field 5 = day of week
\end{itemize}

For \texttt{*/30 * * * *} specifically:
\begin{itemize}
	\item \texttt{*/30} in the minute field means ``every 30 minutes'' (minute 0 and 30)
	\item Each \texttt{*} means ``any value'' for that field
	\item So the schedule runs every hour, every day, every month
\end{itemize}

Additional examples:
\begin{itemize}
	\item \texttt{0 9 * * *}: run every day at 9:00
	\item \texttt{0 */4 * * *}: run every 4 hours on the hour
	\item \texttt{30 17 * * 1-5}: run at 17:30 on weekdays
\end{itemize}

\subsection{Webhooks}

Use webhooks for external event-driven automation.

For detailed relay and webhook configuration, see the README webhook section:
\url{../README.md#discord-integration-relay---gateway}

\subsubsection{Key Fields}
\begin{itemize}
	\item Source ID
	\item Prompt Template
	\item Target Agent and Target Session
	\item Priority and Paused state
	\item Optional secret reference
\end{itemize}

\subsubsection{Suggested Rollout}
\begin{enumerate}
	\item Start with a simple template
	\item Send one synthetic test event
	\item Confirm output in chat/activity panel
	\item Expand template complexity gradually
\end{enumerate}

\subsection{LLM Dependency Note}

Provider/model switching lives in \textbf{General}. Ensure the corresponding provider key exists in API Keys (for example \texttt{OPENAI\_API\_KEY}).

\subsection{Voice}

Use this section to configure wake-word listening, speech capture behavior, and browser text-to-speech playback.

\subsubsection{What You Can Configure}
\begin{itemize}
	\item \textbf{Wake word}: enables or disables the browser-based listener used to detect ``hey proxi''
	\item \textbf{Auto-stop sensitivity}: sets how many seconds of silence are required before Proxi stops recording
	\item \textbf{Auto-send after silence}: controls whether Proxi sends the transcript automatically when recording stops
	\item \textbf{Start-listening beep}: toggles the audible beep when voice capture begins
	\item \textbf{Text to speech}: enables or disables spoken assistant replies
	\item \textbf{Voice / accent}: selects the browser voice used for narration
	\item \textbf{Narration speed}: adjusts how fast spoken replies are read
	\item \textbf{Preview Voice}: plays a short sample with the current voice settings
\end{itemize}

\subsubsection{Suggested Usage}
\begin{itemize}
	\item Turn on wake word if you want hands-free interaction in the GUI
	\item Leave auto-send on for the fastest voice-first workflow
	\item Use Preview Voice after picking a new browser voice or changing speed
	\item Keep TTS off if you only want transcription and do not want spoken replies
\end{itemize}

\subsubsection{Main Screen Voice Controls}
The GUI chat composer also shows the current voice state beside the session info, and the Mic button starts or stops the current recording session.

\subsection{MCP Setup by Integration}

This section explains setup for each requested MCP.

\subsubsection{Google Calendar MCP}
\begin{quote}
\noindent\textbf{Required Keys}
\begin{verbatim}
GMAIL_CLIENT_ID
GMAIL_CLIENT_SECRET
\end{verbatim}

\noindent\textbf{Important Note}\\
Google Calendar uses the same Google OAuth credentials as Gmail in this project.

\noindent\textbf{Setup from Google Cloud Platform (GCP)}
\begin{enumerate}
	\item Open Google Cloud Console and create/select a project
	\item Enable \textbf{Google Calendar API} and \textbf{Gmail API}
	\item Configure OAuth consent screen (External/Internal as needed)
	\item Create OAuth credentials as \textbf{Desktop App}
	\item Copy the generated Client ID and Client Secret into:
	\begin{itemize}
		\item \texttt{GMAIL\_CLIENT\_ID}
		\item \texttt{GMAIL\_CLIENT\_SECRET}
	\end{itemize}
\end{enumerate}

\medskip
\noindent\textbf{First-Time Authorization Behavior}\\
On first use of Google Calendar or Gmail tools, your browser opens Google sign-in and consent pages. After you approve access, Proxi stores a token file for future requests so you are not prompted every time.

\medskip
\noindent\textbf{Validation Prompt}\\
``List my upcoming calendar events for tomorrow.''
\end{quote}

\subsubsection{Gmail MCP}
\begin{quote}
\noindent\textbf{Required Keys}
\begin{verbatim}
GMAIL_CLIENT_ID
GMAIL_CLIENT_SECRET
\end{verbatim}

\noindent\textbf{Setup Note}\\
Gmail uses the same GCP OAuth setup described in \textbf{Google Calendar MCP} above. You only need one Google OAuth client for both integrations.

\medskip
\noindent\textbf{First-Time Authorization Behavior}\\
On first Gmail action, browser-based Google consent is required once to generate/store token credentials.

\medskip
\noindent\textbf{Validation Prompt}\\
``Show my latest unread emails.''
\end{quote}

\subsubsection{Notion MCP}
\begin{quote}
\noindent\textbf{Required Keys}
\begin{verbatim}
NOTION_API_KEY
NOTION_PARENT_PAGE_ID
\end{verbatim}

\noindent\textbf{Where to Get the Notion Values}
\begin{enumerate}
	\item Go to Notion integrations (My Integrations) and create an internal integration
	\item Copy the integration secret into \texttt{NOTION\_API\_KEY}
	\item Open the target parent page in Notion and copy its page ID from the URL
	\item Save that ID as \texttt{NOTION\_PARENT\_PAGE\_ID}
\end{enumerate}

Also share the target page with your integration, otherwise API calls can fail with permissions errors.

\medskip
\noindent\textbf{Validation Prompt}\\
``Create a new Notion note called Weekly Plan.''
\end{quote}

\subsubsection{Obsidian MCP}
\begin{quote}
\noindent\textbf{Required Inputs}\\
This MCP typically uses vault path/configuration rather than cloud API secrets.

\noindent\textbf{Vault Path Setup}\\
Set your Obsidian vault path in configuration/environment used by your Proxi runtime.

Example vault path values:
\begin{itemize}
	\item Windows: \texttt{C:/Users/YOUR\_USERNAME/Documents/ObsidianVault}
	\item macOS: \texttt{/Users/YOUR\_USERNAME/Documents/ObsidianVault}
	\item Linux: \texttt{/home/YOUR\_USERNAME/Documents/ObsidianVault}
\end{itemize}

Ensure the path exists and Proxi has permission to read/write notes there.

\medskip
\noindent\textbf{Validation Prompt}\\
``Create an Obsidian note named Meeting Notes in my vault.''
\end{quote}

\subsubsection{Spotify MCP}
\begin{quote}
\noindent\textbf{Required Keys}
\begin{verbatim}
SPOTIFY_CLIENT_ID
SPOTIFY_CLIENT_SECRET
\end{verbatim}

\noindent\textbf{Setup from Spotify Developer Dashboard}
\begin{enumerate}
	\item Open Spotify Developer Dashboard and create an app
	\item Copy app Client ID into \texttt{SPOTIFY\_CLIENT\_ID}
	\item Copy app Client Secret into \texttt{SPOTIFY\_CLIENT\_SECRET}
	\item Add the redirect URI used by your project (for example \texttt{http://127.0.0.1:8888/callback})
\end{enumerate}

\medskip
\noindent\textbf{First-Time Authorization Behavior}\\
On first Spotify action, a browser window opens for Spotify login/consent. After approval, a token file is generated and reused for subsequent requests.

\medskip
\noindent\textbf{Validation Prompt}\\
``Play my focus playlist on Spotify.''
\end{quote}

\subsubsection{Weather MCP}
\begin{quote}
\noindent\textbf{Required Inputs}\\
No API key is required in this setup. Weather uses the built-in Open-Meteo integration.

\medskip
\noindent\textbf{Setup Note}\\
Enable \textbf{Weather} in \textbf{Settings \textgreater\ Integrations}. Accurate location and timezone settings in your profile help weather responses stay relevant.

\medskip
\noindent\textbf{Validation Prompt}\\
``What is the weather in Hamilton, ON today?''
\end{quote}

\subsubsection{Web MCP}
\begin{quote}
\noindent\textbf{Required Inputs}\\
No API key is required in this setup. Web uses Proxi's built-in web search integration.

\medskip
\noindent\textbf{Setup Note}\\
Enable \textbf{Web} in \textbf{Settings \textgreater\ Integrations} when you want web search and browser-style lookup tools available to the assistant.

\medskip
\noindent\textbf{Validation Prompt}\\
``Search the web for the latest news about OpenAI.''
\end{quote}

\subsection{Example Setup: Complete Key and MCP Flow}

\subsubsection{Step 1: Add Keys}
Add all required key names in \textbf{Settings \textgreater\ API Keys}:
\begin{verbatim}
OPENAI_API_KEY
GMAIL_CLIENT_ID
GMAIL_CLIENT_SECRET
NOTION_API_KEY
NOTION_PARENT_PAGE_ID
DISCORD_BOT_TOKEN
PROXI_DISCORD_COMMAND_PREFIX
SPOTIFY_CLIENT_ID
SPOTIFY_CLIENT_SECRET
\end{verbatim}

Use \texttt{!proxi} as the value for \texttt{PROXI\_DISCORD\_COMMAND\_PREFIX}.

\subsubsection{Step 2: Enable Integrations}
Enable these integrations in \textbf{Settings \textgreater\ Integrations}:
\begin{itemize}
	\item Google Calendar
	\item Gmail
	\item Notion
	\item Obsidian
	\item Spotify
	\item Weather
	\item Web
\end{itemize}

\subsubsection{Step 3: Validate One-by-One}
Run one prompt per MCP to verify setup before moving to automation.

\section{Discord Setup}

This section contains the full Discord relay configuration and usage flow.

\subsection{Configure Discord Source in Gateway}

Add or update a source in your gateway configuration:

\begin{verbatim}
sources:
  discord:
	type: discord
	target_agent: work
	target_session: discord
	priority: 0
	paused: false
	# Optional:
	# discord_command_prefix: /proxi
	# discord_allow_plain: false
	# discord_session_mode: fixed   # fixed | channel | user
	# discord_agent_overrides: {}
\end{verbatim}

Session mode behavior:
\begin{itemize}
	\item \texttt{fixed}: all Discord messages for this source use one session
	\item \texttt{channel}: each Discord channel uses its own session
	\item \texttt{user}: each user in each channel uses its own session
\end{itemize}

\subsection{Security Environment Variable (Recommended)}

In the gateway process environment:

\begin{verbatim}
$env:DISCORD_WEBHOOK_SECRET = "your-shared-secret"
\end{verbatim}

When set, \texttt{/channels/discord/webhook} requires \texttt{X-Signature-256} HMAC signatures.

\subsection{Start Discord Relay}

From project root (manual):

\begin{verbatim}
cd discord_relay
bun install
bun run start
\end{verbatim}

Or via unified CLI:

\begin{verbatim}
uv run proxi discord
\end{verbatim}

\texttt{proxi discord} requires a running gateway and exits if gateway is unreachable.

Relay environment template: \texttt{discord\_relay/.env.example}

\subsection{Discord Command Usage}

Default prefix is \texttt{/proxi}.

\begin{itemize}
	\item \texttt{/proxi <task>}: enqueue a Proxi task
	\item \texttt{/proxi status}: show lane status for this channel/user session
	\item \texttt{/proxi abort}: abort active run in this session
	\item \texttt{/proxi switch <agent\_id>}: route channel to a different agent
	\item \texttt{/proxi help}: show command help
\end{itemize}

To allow plain (non-prefixed) messages, set \texttt{discord\_allow\_plain: true} in gateway source config and \texttt{PROXI\_DISCORD\_ALLOW\_PLAIN=1} in relay environment.

\section{Quick End-to-End Flow}

\begin{enumerate}
	\item Start gateway once: \texttt{uv run proxi gateway start}
	\item Launch TUI (\texttt{uv run proxi}) or GUI (\texttt{uv run proxi frontend})
	\item In GUI settings, configure General and Agent/Profile first
	\item Add API keys and then enable required integrations
	\item Add a cron job using Easy Schedule (Every + Unit), or use manual cron if needed
	\item Send test prompts to confirm integrations and automation are working
\end{enumerate}


\end{document}